\documentclass[11pt, a4paper]{article}

% --- UNIVERSAL PREAMBLE BLOCK ---
\usepackage[a4paper, top=2.5cm, bottom=2.5cm, left=2cm, right=2cm]{geometry}
\usepackage{fontspec}
\usepackage[english, bidi=basic, provide=*]{babel}
\babelprovide[import, onchar=ids fonts]{english}

% Set default/Latin font to Sans Serif in the main (rm) slot
\babelfont{rm}{Noto Sans}

% --- SMART PACKAGE LOADING ---
\usepackage{amsmath}
\usepackage{amssymb}
\usepackage{booktabs}
\usepackage{enumitem}
\usepackage{tabularx}
\usepackage{caption}
\usepackage{microtype}

% hyperref must be the very last package loaded
\usepackage[colorlinks=true, linkcolor=blue, citecolor=blue, urlcolor=blue]{hyperref}

\title{\textbf{In-Depth Literature Review of Bayesian Nonparametric Models for Zero-Inflated Continuous Data}}
\author{\textbf{Literature Review Series in Bayesian Statistics}}
\date{\today}

\begin{document}
\maketitle

\begin{abstract}
Continuous outcomes characterized by an excess of zeros—commonly referred to as semicontinuous or zero-inflated continuous data—present unique challenges in statistical modeling. This literature review synthesizes advanced Bayesian nonparametric and semiparametric frameworks developed to handle these data structures. We explore two-part formulations, variable selection mechanisms under spike-and-slab and adaptive lasso priors, joint generative Dirichlet Process and Enriched Dirichlet Process architectures, unified semicontinuous Tweedie density estimators, and nonparametric copula specifications for zero-inflated margins. Finally, we discuss flexible priors beyond the Dirichlet Process, such as Polya Trees and the Sparsity-Inducing Partition prior, concluding with a comparative methodological synthesis.
\end{abstract}

\section{Theoretical Characterization of Semicontinuous and Zero-Inflated Continuous Data}

In statistical modeling, researchers frequently encounter data characterized by a mixture of a discrete point mass at zero and a continuous distribution of positive values.\cite{olsen2001, tooze2002, albert2005} In health services research, climatology, and financial economics, these data are formally termed semicontinuous or zero-inflated continuous outcomes.\cite{olsen2001, craiu2012, duan1983} Semicontinuous data must be distinguished from zero-inflated count data, such as those modeled by Zero-Inflated Poisson (ZIP) or Zero-Inflated Negative Binomial (ZINB) frameworks.\cite{lambert1992} While count models handle discrete, integer-valued events, semicontinuous modeling deals with nonnegative, continuous outcomes where a substantial subpopulation records exact zeros.\cite{olsen2001, albert2005, duan1983} Examples include individual medical expenditures, where zeros represent patients who do not use health services and positive values describe the skewed continuous expenditure levels among health system users.\cite{olsen2001, neelon2010} Similarly, in nutritional epidemiology, 24-hour dietary recalls for episodically consumed foods include exact zeros for non-consumption days alongside positive continuous measurements for consumption days.\cite{tooze2006, kipnis2009}

Historically, continuous outcomes with excessive zeros were modeled by adding a small positive constant $c$ and performing linear regressions on the log-transformed values, $\log(Y + c)$.\cite{duan1983} This approach, known as the ``one-part'' transformation model, is theoretically problematic.\cite{duan1983} It preserves a point mass at $\log(c)$ on the transformed scale and makes downstream statistical inference highly sensitive to the arbitrary choice of $c$.\cite{duan1983} Additionally, this method fails to distinguish the distinct biological or economic mechanisms that generate the zero state versus the magnitude of the positive continuous response.\cite{tooze2002, duan1983}

To address these limitations, semicontinuous data are represented as a two-part mixture comprising a degenerate distribution at zero and a continuous probability distribution on the positive real line.\cite{lachenbruch2002, min2002, kokonendji2021} Formally, let $Y_i$ represent the observed semicontinuous outcome for subject $i$.\cite{kokonendji2021, muller2004} Semicontinuous outcomes admit the representation:
\begin{equation}
Y_i = \delta_i Z_i
\end{equation}
where $\delta_i$ is a binary indicator representing whether the outcome is positive, $\text{P}(\delta_i = 1) = 1 - p_0$, and $Z_i$ is a positive continuous random variable with a distribution function $F$ supported on $(0, \infty)$ and a density $f$ with respect to the Lebesgue measure.\cite{albert2005, kokonendji2021, muller2004} In compositional data analysis (CoDA), this structure relates directly to the ``compositional zero problem''.\cite{aitchison1982, aitchison1986} Compositional data are defined on the simplex, where the sum of all components is constrained to a constant, such as one.\cite{aitchison1982, aitchison1986} Traditional log-ratio transformations, including the Additive Log-Ratio (ALR) and Isometric Log-Ratio (ILR), map compositions from the simplex to real space but cannot directly accommodate zeros because logarithmic operations require strictly positive arguments.\cite{aitchison1986} Standard imputation approaches that replace zeros with small pseudo-counts distort the covariance structure and bias downstream analyses.\cite{aitchison1982, aitchison1986, egozcue2003}

To preserve the structural meaning of zeros in multivariate compositional continuous spaces, researchers developed zero-inflated log-normal graphical models (Zi-LN).\cite{fang2015, kurtz2015} These models relate observed proportions to a latent multivariate normal variable $z_i$ and a physical abundance variable $a_i$.\cite{fang2015, kurtz2015} The continuous abundance follows a zero-inflated multivariate distribution:
\begin{equation}
a_i \sim \mathcal{N}(\mu, \Sigma) \mathbb{I}(z_i > \delta)
\end{equation}
where $\mathbb{I}(\cdot)$ is an indicator function, $\delta$ determines the zero-inflation threshold, and the precision matrix $\Sigma^{-1}$ encodes conditional dependencies in a Gaussian graphical framework.\cite{fang2015, kurtz2015} This formulation allows ``biological'' zeros to be represented as structurally absent states governed by a latent continuous process, rather than treating them as arbitrary sampling zeros.\cite{kurtz2015}


\section{Semiparametric Two-Part Formulations and Variable Selection Mechanisms}

Two-part regression models are the standard framework for analyzing semicontinuous data in longitudinal and spatial settings.\cite{olsen2001, tooze2002, neelon2010} A typical parametric two-part model uses a Lognormal (LN) density for the positive component, modeling the probability of a positive response via a binary link function (such as probit or logit) and the mean of the log-scale positive expenditures via linear regression.\cite{olsen2001, neelon2010} However, lognormal assumptions impose strict symmetry and unimodality on the log scale, which is often violated in empirical applications characterized by multi-modal, heavy-tailed, or heteroscedastic positive distributions.\cite{olsen2001, manning2001} 

To relax these parametric assumptions, researchers proposed semiparametric Bayesian two-part models.\cite{olsen2001, neelon2010, muller2004} Under these specifications, the binary selection indicator $\delta_i$ is modeled via a semiparametric Bernoulli regression, while the joint distribution of the positive continuous response and continuous predictors is modeled using a Dirichlet Process Mixture (DPM) of Normals.\cite{neelon2010, muller2004, manning2001} The DPM prior allows the continuous portion of the hurdle model to adaptively capture skewness and locate latent subject clusters, providing robust predictions without requiring the pre-specification of higher-order interaction terms.\cite{neelon2010, muller2004}

In high-dimensional settings with many candidate predictors, determining which explanatory variables influence the zero and continuous components is a key challenge.\cite{min2002, george1993, ishwaran2005} Statisticians addressed this by developing a semiparametric Bernoulli-Normal Two-part (BNT) regression model with Bayesian variable selection.\cite{min2002, george1993, ishwaran2005} For a semicontinuous outcome $s_i$, the model identifies the point mass at zero using two surrogate variables:
\begin{equation}
u_i = \mathbb{I}(s_i > 0) \quad \text{and} \quad z_i = \log(s_i + \epsilon)
\end{equation}
where $\epsilon$ is a positive constant that allows the log transformation of zero values.\cite{ishwaran2005} The BNT model utilizes two distinct regression components to capture the selection and continuous pathways:
\begin{align}
\text{Part 1: } \quad & \text{P}(u_i = 1 \mid X_i) = \Phi(X'_i \beta_1) \\
\text{Part 2: } \quad & z_i \mid u_i = 1, X_i \sim \mathcal{N}(X'_i \beta_2, \sigma^2)
\end{align}
where $\Phi(\cdot)$ denotes the standard normal cumulative distribution function, and $X_i$ is a vector of observed covariates.\cite{min2002, ishwaran2005} 

To select active predictors from a large pool of covariates, researchers specify Bayesian bimodal spike-and-slab bimodal priors or Bayesian Lasso priors on the regression coefficients $\beta_1$ and $\beta_2$.\cite{george1993, ishwaran2005} The spike-and-slab bimodal prior specifies a mixture distribution for each coefficient:
\begin{equation}
\beta_j \sim (1 - \gamma_j) \mathcal{N}(0, \sigma_0^2) + \gamma_j \mathcal{N}(0, \sigma_1^2)
\end{equation}
where the variance parameter $\sigma_0^2$ is extremely small (the ``spike''), forcing the coefficient toward zero, and $\sigma_1^2$ is large (the ``slab''), allowing the model to estimate a non-zero effect.\cite{george1993, ishwaran2005} The binary indicator $\gamma_j \sim \text{Bernoulli}(w_j)$ determines the posterior probability of inclusion for each covariate.\cite{george1993, ishwaran2005} Under this hierarchical construction, the model adaptively adjusts the level of shrinkage across both parts of the regression, identifying key predictors and unobserved latent factors.\cite{george1993, ishwaran2005}

For longitudinal studies featuring repeated measures, two-part models must also account for within-subject correlation and temporal trends.\cite{tooze2002} For instance, in clinical substance use research, primary outcomes such as alcohol or drug consumption are semicontinuous.\cite{neelon2011} A three-level Bayesian mixed-effects semicontinuous model was developed to analyze such repeated-measure datasets:
\begin{align}
\text{Level 3: } \quad & \xi \sim p(\xi) \\
\text{Level 2: } \quad & \phi = (\phi_0^\top, \phi_c^\top)^\top \sim p(\phi \mid \xi) \\
\text{Level 1 (latent): } \quad & w_{i,0} \sim \mathcal{N}(x_{i,0}\beta_0 + M_i\phi_0, 1), \quad u_i = \mathbb{I}(w_{i,0} > 0) \\
& y_i^* \sim \mathcal{N}(x_{i,c}\beta_c + M_i\phi_c, \sigma_y^2) \\
\text{Level 1 (observed): } \quad & y_i = u_i y_i^*
\end{align}
where $\beta_0$ and $\beta_c$ are fixed-effects vectors, $M_i$ is a design matrix, and $\phi = (\phi_0^\top, \phi_c^\top)^\top$ is a vector of subject-specific random effects capturing correlation between the selection and continuous components.\cite{neelon2011} To relax the assumption of normality for the random effects, which can obscure subgroup structures, researchers specify a DPM of Normals as the prior for $\phi$, enabling the model to nonparametrically identify latent patient trajectories and subgroups.\cite{neelon2010}


\section{Advanced Dirichlet Process Architectures and Generative Joint Models}

The development of advanced Bayesian nonparametric architectures has produced highly flexible models for predicting semicontinuous outcomes and estimating causal treatment effects.\cite{oganisian2021, oganisian2020} A prominent contribution is the multipurpose joint model proposed by Oganisian, Mitra, and Roy (2020), which is implemented via the Chinese Restaurant Process (CRP) in the \texttt{ChiRP} package.\cite{oganisian2020, oganisianpkg, neal2000} Unlike standard conditional models that treat the covariates $X_i$ as fixed and model only the conditional density $Y_i \mid X_i$, this approach specifies a joint generative model for the outcome $Y_i$ and the covariates $X_i$.\cite{oganisian2020}

Let $Y_i$ be the zero-inflated continuous outcome, $A_i$ represent the treatment indicator, and $L_i$ denote a vector of observed baseline confounders, so that the full covariate vector is $X_i = (1, A_i, L_i)'$.\cite{oganisian2020} The joint distribution of the outcome and the covariates is modeled as a mixture of localized parametric components governed by a Dirichlet Process prior.\cite{oganisian2020} Let $\omega_i = (\beta_i, \phi_i, \gamma_i, \eta_i, \theta_i)$ represent the cluster-specific parameters for subject $i$.\cite{oganisian2020} The conditional distribution of the continuous outcome $Y_i$ is modeled as a two-part mixture:
\begin{align}
Y_i \mid X_i, \omega_i & \sim \pi(X'_i \gamma_i) \delta_0(y_i) + \left[1 - \pi(X'_i \gamma_i)\right] \mathcal{N}(X'_i \beta_i, \phi_i) \\
\omega_i \mid G & \sim G \\
G & \sim \text{DP}(\alpha, G_0)
\end{align}
where $\pi(\cdot)$ represents a probit link function, $\delta_0(y_i)$ represents a point mass at zero, and the marginal distribution of the covariates is modeled nonparametrically via localized parameters $\eta_i$ and $\theta_i$.\cite{oganisian2020} By estimating the full joint distribution, this generative formulation handles missing-at-random (MAR) covariates through data augmentation, generates highly flexible predictive distributions, and induces a natural clustering of subjects with similar clinical profiles.\cite{oganisian2020, oganisianpkg}

However, when modeling a large number of baseline covariates, standard joint DP mixtures suffer from a structural limitation.\cite{wade2011} Because a single DP governs the partitioning of both the outcome and the covariates, the high-dimensional covariate space can dominate the clustering process, placing subjects into clusters based on their covariates rather than their outcome patterns.\cite{wade2011} This reduces the model's sensitivity to the outcome distribution and degrades predictive accuracy.\cite{wade2011}

To decouple the clustering of covariates from the outcome model, researchers proposed the Enriched Dirichlet Process (EDP) mixture model.\cite{oganisian2021, wade2011, wade2014} The EDP defines a nested clustering structure, denoted by the cluster membership $C_i = (C_{i, y}, C_{i, x|y})$ for observation $i$.\cite{wade2011} The outer cluster membership $C_{i, y}$ is driven by the outcome model structure ($Y_i \mid A_i, X_i$), while the nested cluster membership $C_{i, x|y}$ is driven by the baseline covariates $X_i$ nested within the outer outcome cluster.\cite{wade2011} For each sampled cluster, the model updates the outcome parameters $\phi_c$ and the nested covariate parameters $\gamma_{j|c}$ from their respective conditional posteriors utilizing base measures $G_{0, \phi}$ and $G_{0, \gamma}$ derived from the base measure $G_0$.\cite{wade2011}

Using the posterior samples of the parameters, the Conditional Average Treatment Effect (CATE; $\tau(x)$) and the zero-inflation probability are calculated using conditional expectations that incorporate the nested cluster weights $\omega_c(l_i)$, where $l_i$ denotes the individual covariate profile:
\begin{align}
\text{E}(Y_i \mid A_i = a, X_i = x, \phi, \gamma, C) = & \, \frac{\alpha_\phi}{\alpha_\phi + n} \int K(l_i \mid \gamma) dG_{0, \gamma}(\gamma) \int \text{E}(Y \mid l_i, \phi) dG_{0, \phi}(\phi) \nonumber \\
& + \sum_{c=1}^{C_y} \omega_c(l_i) \text{E}(Y \mid l_i, \phi_c) \\[10pt]
\text{P}(Y_i = 0 \mid A_i = a, X_i = x, \phi, \gamma, C) = & \, \frac{\alpha_\phi}{\alpha_\phi + n} \int K(l_i \mid \gamma) dG_{0, \gamma}(\gamma) \int \text{P}(Y = 0 \mid l_i, \phi) dG_{0, \phi}(\phi) \nonumber \\
& + \sum_{c=1}^{C_y} \omega_c(l_i) \text{P}(Y = 0 \mid l_i, \phi_c)
\end{align}
This nested clustering structure allows the model to capture pre-treatment heterogeneity and trace back subgroups with significant treatment effects.\cite{oganisian2021, wade2011} 

In causal inference settings, this EDP model exhibits distinct advantages over tree-based methods such as Bayesian Additive Regression Trees (BART) or Bayesian Causal Forests (BCF).\cite{oganisian2021, wade2011} Tree-based methods are sensitive to minor data variations and struggle to accurately reflect posterior uncertainty in regions where the covariate overlap (positivity) assumption is violated.\cite{oganisian2021, wade2011} In contrast, the EDP-based model naturally assigns subjects to sparse covariate clusters.\cite{wade2011} In non-overlap regions of the covariate space, the model propagates this lack of data density directly into the posterior, generating wider credible intervals that accurately reflect the underlying scientific uncertainty.\cite{oganisian2021, wade2011}

An alternative approach is the Enriched Finite Mixture model with a random number of components, which generalizes the finite mixture of hurdle models.\cite{dellaportas2006, richardson1997, papageorgiou2011, green1995} This framework models multiple zero-inflated continuous processes by specifying a hurdle model for each process.\cite{dellaportas2006, papageorgiou2011, green1995} Conditional on the model parameters, the different processes are assumed independent, which reduces the number of parameters compared to traditional multivariate approaches.\cite{dellaportas2006, papageorgiou2011} 

The subject-specific probabilities of zero-inflation and the parameters of the sampling distribution are modeled via an enriched finite mixture with a random number of components.\cite{dellaportas2006, papageorgiou2011} This induces a two-level clustering of the subjects based on the patterns of zero/non-zero values (outer clustering) and the continuous positive sampling distribution (inner clustering).\cite{dellaportas2006, richardson1997, papageorgiou2011} The number of components is treated as a random variable to allow for a fully Bayesian treatment of model complexity.\cite{papageorgiou2011, miller2018}


\section{Unified Semicontinuous Tweedie and Nonparametric Density Estimators}

An alternative to two-part specifications is the family of Tweedie exponential dispersion models.\cite{albert2005, kokonendji2021, tweedie1984} Semicontinuous outcomes are naturally accommodated by a Tweedie distribution when the power index $p$ lies strictly between $1$ and $2$.\cite{albert2005, kokonendji2021, tweedie1984} Under this parameterization, the Tweedie distribution behaves as a compound Poisson-Gamma process, representing the sum of a Poisson number of independent Gamma-distributed random variables.\cite{tweedie1984} The continuous probability density function for an outcome $y \ge 0$ is formulated as:
\begin{equation}
f(y; \mu, \phi, p) = a(y, \phi, p) \exp \left\{ \frac{1}{\phi} \left[ y \frac{\mu^{1-p}}{1-p} - \frac{\mu^{2-p}}{2-p} \right] \right\}
\end{equation}
where $\mu > 0$ is the mean, $\phi > 0$ is the dispersion parameter, and $a(y, \phi, p)$ is a normalizing constant.\cite{tweedie1984} The discrete point mass at zero arises intrinsically from the Poisson zero-state of the compound process and is written as:
\begin{equation}
\text{P}(Y = 0) = \exp(-\lambda), \quad \text{where } \lambda = \frac{\mu^{2-p}}{\phi (2-p)}
\end{equation}
In classical settings, Zero-Inflated Tweedie (ZIT) models extend this structure by incorporating an explicit Bernoulli mixing variable to handle extremely unbalanced nonnegative data.\cite{tweedie1984} 

From a nonparametric perspective, traditional two-part density estimators are fundamentally asymmetric, smoothing the positive component on $(0, \infty)$ using only the positive observations and treating the zero-inflation probability $p_0$ entirely separately.\cite{albert2005, kokonendji2021} This separation causes severe boundary bias near zero and makes it extremely difficult to recover the true shape of the continuous density in close proximity to the point mass.\cite{albert2005, kokonendji2021} 

To overcome this, researchers developed a Tweedie-based nonparametric kernel estimator for the full semicontinuous density.\cite{albert2005, kokonendji2021} Because the Tweedie kernel (for $1 < p < 2$) possesses its own native point mass at zero and an absolutely continuous component on the positive real line, it yields a unified smoothing construction.\cite{kokonendji2021} Rather than adapting a continuous kernel to the positive half-line, the Tweedie kernel estimator uses the \textit{entire} semicontinuous sample simultaneously.\cite{kokonendji2021} 

At the boundary $x = 0$, the estimator coincides exactly with the empirical proportion of zeros, while for $x > 0$, the zero-valued observations continue to constructively influence the estimation of the continuous component near the boundary.\cite{kokonendji2021} This boundary-aware, unified smoothing approach drastically reduces boundary bias and significantly improves the estimation of heavy-tailed, highly skewed semicontinuous densities.\cite{albert2005, kokonendji2021}


\section{Nonparametric Copula Specifications for Zero-Inflated Margins}

When analyzing multivariate semicontinuous data, researchers must model both the marginal distributions and the complex, potentially non-linear dependence structures between multiple zero-inflated variables.\cite{craiu2012, tooze2006, schepsmeier2016} Copula modeling, anchored in Sklar's theorem, provides a natural framework for this task by decomposing a joint distribution into its marginals and a copula function $C: [0, 1]^d \to [0, 1]$ that fully describes the dependence structure.\cite{craiu2012, sklar1959, joe2014}

However, the application of copulas to semicontinuous data is hindered by a major theoretical hurdle: Sklar's theorem only guarantees a unique copula representation when the marginal distributions are strictly continuous.\cite{craiu2012, joe2014, genest2007} For mixed or semicontinuous data, which feature massive ties at zero, the probability integral transform $F_j(Y_j)$ is not uniformly distributed, rendering standard nonparametric copula estimators inconsistent and parametric specifications highly unstable.\cite{craiu2012, joe2014, genest2007} 

To resolve this issue, statisticians developed consistent nonparametric copula estimators tailored for mixed and semicontinuous outcomes.\cite{craiu2012} These estimators utilize local averaging approaches and continuous covariates to expand the support of the marginals, restoring identifiability and enabling consistent estimation of the ``hidden'' dependence structure.\cite{craiu2012, joe2014, genest2007} 

In high-dimensional climatological and biological applications, this methodology was further advanced through Vine Copula Bias Correction (VBC).\cite{schepsmeier2016} This technique develops a new vine density decomposition specifically designed to accommodate semicontinuous variables.\cite{schepsmeier2016} To handle the discrete-continuous mixture at zero, the model employs a randomized version of the Rosenblatt transform, which projects the zero-inflated variables into a uniform probability space without introducing artificial bias.\cite{schepsmeier2016} This randomized transformation allows the joint dependence structure to be adjusted and transferred consistently across domains.\cite{schepsmeier2016}

This copula framework is widely implemented in hierarchical latent variable models.\cite{tooze2006} For example, in nutritional epidemiology, long-term dietary intake of episodically consumed foods must be estimated from 24-hour recalls that contain exact zeros.\cite{tooze2006} Researchers address this by designing a copula-based approach that translates the semicontinuous observed recalls into continuous latent surrogates, modeling the joint distribution using flexible nonparametric marginal specifications while preserving the exact zero-consumption patterns in the observed data.\cite{tooze2006}


\section{Nonparametric Priors Beyond the Dirichlet Process: Polya Trees and Generalizations}

While Dirichlet Process mixtures represent the dominant BNP paradigm, alternative nonparametric priors offer distinct theoretical and computational advantages for modeling continuous zero-inflated outcomes.\cite{miller2018, barrientos2012, lavine1992} Chief among these is the Polya Tree (PT) prior.\cite{lavine1992} 

Polya trees construct random probability measures by recursively partitioning the sample space and assigning mass to child elements at each partition level according to Beta distributions.\cite{lavine1992} A fundamental limitation of the Dirichlet Process is that a random measure drawn from a DP is almost surely discrete.\cite{lavine1992} To obtain a continuous density from a DP, one must convolve it with a continuous kernel, forming a DPM model.\cite{lavine1992} In contrast, Polya tree priors can assign probability 1 directly to the class of absolutely continuous distributions, bypassing the need for kernel mixtures.\cite{egozcue2003}

In zero-inflated settings, Polya trees are increasingly utilized in dynamic model selection and dynamic association frameworks.\cite{hanson2002, hanson2006} For instance, a zero-inflated conditionally identically distributed (CID) species sampling prior allows observations to be clustered into either a null parametric process or a flexible nonparametric Polya tree process, permitting spatial and temporal variation in model selection.\cite{hanson2006} 

Furthermore, researchers developed hierarchical and stateful Polya trees, which employ latent discrete state variables at each partition branch to model local smoothness and concentration parameters nonparametrically across different regions of the sample space.\cite{lavine1992} This local adaptability is highly beneficial for semicontinuous data, where the density may exhibit an extremely sharp peak near zero but a smooth, heavily skewed tail elsewhere.\cite{lavine1992}

Beyond Polya trees, other innovative prior formulations have been introduced to refine weight clustering and sparsity in nonparametric mixtures.\cite{miller2018, barrientos2012, ishwaran2001} The Enriched Pitman-Yor (EPY) process generalizes the enriched Dirichlet process, providing a unified probabilistic framework that accommodates DP mixtures with spike-and-slab base measures and mixture-of-mixtures formulations while preserving computational tractability.\cite{barrientos2012, ishwaran2001} 

Additionally, to prevent the over-allocation of small, redundant clusters typical of standard Dirichlet priors, researchers proposed the Sparsity-Inducing Partition (SIP) prior.\cite{miller2018} The SIP prior utilizes the Selberg Dirichlet distribution to govern mixture weights.\cite{miller2018} This distribution incorporates a repulsive term that naturally penalizes component locations that lie close to one another on the simplex, encouraging a sparse partition with few dominant, well-separated clusters, which significantly enhances the interpretability of subject clustering in semicontinuous data.\cite{miller2018}


\section{Comparative Analysis and Methodological Synthesis}

To systematically evaluate the various Bayesian nonparametric and semiparametric frameworks for zero-inflated continuous and semicontinuous data, this section synthesizes the core mathematical properties and variable selection mechanisms. The comparative evaluations focus on structural differences, clustering features, and treatment-effect estimation capabilities.

Table~\ref{tab:frameworks} evaluates the key properties of the main estimator classes discussed in this literature review.

\begin{table}[htbp]
\centering
\caption{Comparative synthesis of Bayesian nonparametric and semiparametric frameworks.}
\label{tab:frameworks}
\small
\begin{tabularx}{\textwidth}{p{3.2cm} X X X p{2.8cm}}
\toprule
\textbf{Estimator Class} & \textbf{Mathematical Representation} & \textbf{Treatment of Zeros} & \textbf{Clustering Mechanism} & \textbf{Implementation Status} \\
\midrule
\textbf{DPM Two-Part (Hurdle)} \cite{olsen2001, muller2004, oganisian2020} & 
$Y_i = \delta_i Z_i$, $Z_i \sim \int \mathcal{N}(\mu, \sigma^2) dG$ & 
Separate binary regression hurdle component.\cite{olsen2001, oganisian2020} & 
Jointly clusters outcomes and covariates via DP.\cite{oganisian2020} & 
Implemented in R package \texttt{ChiRP} (\texttt{ZDPMix}).\cite{oganisianpkg} \\
\addlinespace
\textbf{Enriched DP (EDP)} \cite{oganisian2021, wade2011} & 
Nested partitioning with conditional priors on $G_0$ & 
Multi-component mixture mapping the selection probability.\cite{wade2011} & 
Two-level nested clustering of outcome and covariates.\cite{wade2011} & 
Custom Gibbs sampling using conjugate updates.\cite{wade2011} \\
\addlinespace
\textbf{Semicontinuous Tweedie} \cite{albert2005, kokonendji2021} & 
Compound Poisson-Gamma distribution for $1 < p < 2$ & 
Unified density includes a point mass at zero.\cite{kokonendji2021, tweedie1984} & 
Kernel smoothing across the entire continuous sample.\cite{kokonendji2021} & 
Custom kernel estimators and dispersion models.\cite{kokonendji2021} \\
\addlinespace
\textbf{Nonparametric Copula} \cite{craiu2012, schepsmeier2016} & 
$F(y_1, \dots, y_d) = C(F_1(y_1), \dots, F_d(y_d))$ & 
Randomized Rosenblatt transform of zero-inflated margins.\cite{schepsmeier2016} & 
Local averaging models marginal dependencies.\cite{craiu2012} & 
Custom Vine Copula algorithms (VBC).\cite{schepsmeier2016} \\
\addlinespace
\textbf{Polya Tree (PT)} \cite{lavine1992} & 
Recursive binary partition tree with Beta variables & 
Combined with outer spike-and-slab mixtures.\cite{hanson2006} & 
Localized partitioning of sample spaces.\cite{lavine1992} & 
Custom hierarchical adaptive partition samplers.\cite{lavine1992} \\
\bottomrule
\end{tabularx}
\end{table}

Table~\ref{tab:regularization} reviews the available mechanisms for incorporating high-dimensional covariates and executing model selection, which is a major concern when analyzing sparse healthcare or genomic databases.

\begin{table}[htbp]
\centering
\caption{Regularization methods and variable selection mechanisms in high-dimensional settings.}
\label{tab:regularization}
\small
\begin{tabularx}{\textwidth}{p{3.2cm} X X p{2.5cm} X}
\toprule
\textbf{Regularization Method} & \textbf{Prior Structure} & \textbf{Variable Selection Mechanism} & \textbf{Continuous Support Compatibility} & \textbf{Typical Application} \\
\midrule
\textbf{BART Priors} \cite{chipman2010, linero2018} & 
Sum-of-weak-learners tree structures & 
Splitting frequency determines variable relative importance.\cite{bleich2014} & 
Highly compatible with continuous zero-inflated predictors.\cite{chipman2010, linero2018} & 
Compositional zero probability modeling.\cite{chipman2010, linero2018} \\
\addlinespace
\textbf{Spike-and-Slab DPM} \cite{george1993, ishwaran2005, ishwaranrao2005} & 
Two-point mixture: $\sigma_0^2$ (spike) vs $\sigma_1^2$ (slab) & 
Latent indicators $\gamma_j$ enforce sparse coefficient selection.\cite{george1993} & 
Designed for BNT two-part latent variables.\cite{ishwaran2005} & 
Medical cost regressions and microbiome association studies.\cite{george1993, ishwaranrao2005} \\
\addlinespace
\textbf{Adaptive Lasso BNT} \cite{min2002, ishwaran2005} & 
Laplacians with adaptive penalty weights & 
Shrinks non-informative coefficients to exactly zero.\cite{min2002} & 
Handles log-scale positive outcomes.\cite{ishwaran2005} & 
Dietary intake predictors and economic expenditure surveys.\cite{min2002, ishwaran2005} \\
\addlinespace
\textbf{Enriched Finite Mixtures} \cite{dellaportas2006, papageorgiou2011} & 
Dirichlet priors over nested partition parameters & 
Outer/inner cluster identification reduces parameter space.\cite{dellaportas2006} & 
Designed for shifted multivariate continuous components.\cite{papageorgiou2011} & 
Multivariate clinical processes and questionnaires.\cite{dellaportas2006} \\
\bottomrule
\end{tabularx}
\end{table}


\section{Conclusions}

The statistical modeling of zero-inflated continuous and semicontinuous data has advanced significantly, moving away from restrictive parametric assumptions and arbitrary offset transformations toward robust Bayesian nonparametric (BNP) formulations.\cite{olsen2001, albert2005, duan1983, oganisian2020} Traditional hurdles and two-part models, while useful for separating the zero-generating process from the continuous intensity component, are prone to severe bias under distributional misspecification.\cite{olsen2001, manning2001} The integration of Dirichlet Process Mixtures (DPM) and Enriched Dirichlet Processes (EDP) has resolved these limitations, enabling researchers to capture heavy right-skewness, heteroscedasticity, and multimodality in positive outcomes without sacrificing interpretability or model stability.\cite{olsen2001, oganisian2020, wade2011}

The Enriched Dirichlet Process (EDP) represents a key methodological development, addressing the curse of dimensionality in high-dimensional covariate spaces.\cite{wade2011} By decoupling the outcome clustering from the covariate partition, the EDP prevents covariate noise from destabilizing the outcome model.\cite{wade2011} This structural separation makes the EDP highly effective for estimating heterogeneous treatment effects and capturing posterior uncertainty in regions where the overlap assumption is violated.\cite{oganisian2021, wade2011}

Simultaneously, alternative frameworks have challenged the necessity of modeling the zero and positive components separately.\cite{kokonendji2021} Nonparametric density estimators utilizing Tweedie kernels provide a unified smoothing framework, using the entire sample to reduce boundary bias and accurately reconstruct continuous densities near zero.\cite{albert2005, kokonendji2021} For multivariate semicontinuous data, nonparametric copulas combined with randomized Rosenblatt transformations have resolved the non-uniqueness challenges of Sklar's theorem, providing consistent framework options for high-dimensional, zero-inflated datasets.\cite{craiu2012, schepsmeier2016} 

Future research in this domain lies in the unification of these paradigms---such as combining the nested clustering properties of the Enriched Pitman-Yor process with unified Tweedie-kernel margins or copula-based dependence structures---providing robust, data-adaptive tools for complex, zero-inflated continuous phenomena.


% --- BIBLIOGRAPHY ---
\begin{thebibliography}{47}

\bibitem{olsen2001}
M.~A. Olsen and J.~L. Schafer.
\newblock A two-part random-effects model for semicontinuous longitudinal data.
\newblock {\em Journal of the American Statistical Association}, 96(454):730--745, 2001.

\bibitem{tooze2002}
J.~A. Tooze, G.~K. Grunwald, and R.~H. Jones.
\newblock Analysis of repeated measures semicontinuous data.
\newblock {\em Methods of Information in Medicine}, 41(01):72--77, 2002.

\bibitem{albert2005}
P.~S. Albert and J.~Shen.
\newblock Modeling longitudinal semicontinuous data with an application to an acupuncture clinical trial.
\newblock {\em Biometrics}, 61(4):1111--1120, 2005.

\bibitem{craiu2012}
R.~V. Craiu and A.~Sabeti.
\newblock Inference for copula models with mixed continuous and discrete outcomes.
\newblock {\em Journal of Multivariate Analysis}, 109:251--264, 2012.

\bibitem{duan1983}
N.~Duan, W.~G. Manning, C.~N. Morris, and J.~P. Newhouse.
\newblock A comparison of alternative models for the analysis of medical expenditure data.
\newblock {\em Journal of Business \& Economic Statistics}, 1(2):115--126, 1983.

\bibitem{lambert1992}
D.~Lambert.
\newblock Zero-inflated Poisson regression, with an application to defects in manufacturing.
\newblock {\em Technometrics}, 34(1):1--14, 1992.

\bibitem{neelon2010}
B.~H. Neelon, A.~J. O'Malley, and S.~L. Normand.
\newblock A Bayesian two-part model for longitudinal semicontinuous data.
\newblock {\em Biometrics}, 66(3):856--865, 2010.

\bibitem{tooze2006}
J.~A. Tooze, V.~Kipnis, D.~W. Buckman, R.~J. Carroll, L.~S. Freedman, P.~M. Guenther, and K.~W. Dodd.
\newblock A mixed-effects model for occasional dietary intake.
\newblock {\em Journal of the American Statistical Association}, 101(475):1014--1024, 2006.

\bibitem{kipnis2009}
V.~Kipnis, D.~Midthune, D.~W. Buckman, K.~W. Dodd, P.~M. Guenther, S.~M. Krebs-Smith, and J.~A. Tooze.
\newblock Modeling data with excess zeros and asymmetrical distributions in dietary assessment.
\newblock {\em Statistics in Medicine}, 28(28):3500--3516, 2009.

\bibitem{lachenbruch2002}
P.~A. Lachenbruch.
\newblock Analysis of data with a bunch of zeroes.
\newblock {\em Statistical Methods in Medical Research}, 11(4):297--302, 2002.

\bibitem{min2002}
Y.~Min and A.~Agresti.
\newblock Modeling nonnegative continuous data with a mass at zero.
\newblock {\em Metron}, 60(3--4):1--22, 2002.

\bibitem{kokonendji2021}
C.~C. Kokonendji and S.~M. Som{\'e}.
\newblock Tweedie-based nonparametric kernel density estimation for semicontinuous data.
\newblock {\em Journal of Nonparametric Statistics}, 33(2):263--281, 2021.

\bibitem{muller2004}
P.~M{\"u}ller and F.~A. Quintana.
\newblock Nonparametric Bayesian data analysis.
\newblock {\em Statistical Science}, 19(1):95--110, 2004.

\bibitem{aitchison1982}
J.~Aitchison.
\newblock The statistical analysis of compositional data.
\newblock {\em Journal of the Royal Statistical Society: Series B (Methodological)}, 44(2):139--160, 1982.

\bibitem{aitchison1986}
J.~Aitchison.
\newblock {\em The Statistical Analysis of Compositional Data}.
\newblock Chapman and Hall, London, 1986.

\bibitem{egozcue2003}
J.~J. Egozcue, V.~Pawlowsky-Glahn, G.~Mateu-Figueras, and C.~Barcel{\'o}-Vidal.
\newblock Isometric logratio transformations for compositional data analysis.
\newblock {\em Mathematical Geology}, 35(3):279--300, 2003.

\bibitem{fang2015}
H.~Fang, C.~Huang, H.~Zhao, and M.~Deng.
\newblock CCLasso: correlation inference for compositional data through Lasso.
\newblock {\em Bioinformatics}, 31(19):3111--3118, 2015.

\bibitem{kurtz2015}
J.~D. Kurtz, C.~L. M{\"u}ller, E.~R. Miraldi, L.~J. McIver, H.~Deng, and R.~Bonneau.
\newblock Sparse and compositionally robust inference of microbial ecological networks.
\newblock {\em PLoS Computational Biology}, 11(5):e1004226, 2015.

\bibitem{manning2001}
W.~G. Manning and J.~Mullahy.
\newblock Estimating log-models: to transform or not to transform?.
\newblock {\em Journal of Health Economics}, 20(4):461--494, 2001.

\bibitem{george1993}
E.~I. George and R.~E. McCulloch.
\newblock Variable selection via Gibbs sampling.
\newblock {\em Journal of the American Statistical Association}, 88(423):881--889, 1993.

\bibitem{ishwaran2005}
H.~Ishwaran and J.~S. Rao.
\newblock Spike and slab variable selection: frequentist and Bayesian strategies.
\newblock {\em The Annals of Statistics}, 33(2):730--773, 2005.

\bibitem{neelon2011}
B.~Neelon, B.~Muth{\'e}n, and S.~K. Ghosh.
\newblock A Bayesian mixed-effects semicontinuous model for clinical trials of substance abuse.
\newblock {\em Contemporary Clinical Trials}, 32(3):443--452, 2011.

\bibitem{oganisian2021}
A.~Oganisian, N.~Mitra, and J.~A. Roy.
\newblock A Bayesian nonparametric model for heterogeneous treatment effects of semicontinuous outcomes.
\newblock {\em Biostatistics}, 22(4):794--811, 2021.

\bibitem{oganisian2020}
A.~Oganisian, N.~Mitra, and J.~A. Roy.
\newblock A Bayesian nonparametric joint model for semicontinuous outcomes.
\newblock {\em Biometrics}, 76(4):1105--1116, 2020.

\bibitem{oganisianpkg}
A.~Oganisian.
\newblock {\em ChiRP: Chinese Restaurant Process Mixture Models in R}.
\newblock R package version 0.1.0, 2020.

\bibitem{neal2000}
R.~M. Neal.
\newblock Markov chain Monte Carlo methods for Dirichlet process mixture models.
\newblock {\em Journal of Computational and Graphical Statistics}, 9(2):249--265, 2000.

\bibitem{wade2011}
S.~Wade, S.~Mongelluzzo, and S.~Petrone.
\newblock An enriched Dirichlet process for covariate-dependent clustering.
\newblock {\em Bayesian Analysis}, 6(3):359--385, 2011.

\bibitem{wade2014}
S.~Wade, D.~B. Dunson, S.~Petrone, and L.~Trippa.
\newblock Improving prediction in regression mixture models with the enriched Dirichlet process.
\newblock {\em Journal of Machine Learning Research}, 15(1):1041--1066, 2014.

\bibitem{dellaportas2006}
P.~Dellaportas and I.~Papageorgiou.
\newblock Multivariate mixture of hurdle models.
\newblock {\em Statistics and Computing}, 16(2):163--174, 2006.

\bibitem{richardson1997}
S.~Richardson and P.~J. Green.
\newblock On Bayesian analysis of mixtures with an unknown number of components (with discussion).
\newblock {\em Journal of the Royal Statistical Society: Series B (Statistical Methodology)}, 59(4):731--792, 1997.

\bibitem{papageorgiou2011}
I.~Papageorgiou.
\newblock Enriched finite mixture of hurdle models with a random number of components.
\newblock {\em Computational Statistics \& Data Analysis}, 55(4):1611--1624, 2011.

\bibitem{green1995}
P.~J. Green.
\newblock Reversible jump Markov chain Monte Carlo computation and Bayesian model determination.
\newblock {\em Biometrika}, 82(4):711--732, 1995.

\bibitem{miller2018}
J.~W. Miller and M.~T. Harrison.
\newblock Mixture models with a prior on the number of components.
\newblock {\em Journal of the American Statistical Association}, 113(521):340--356, 2018.

\bibitem{tweedie1984}
M.~C.~K. Tweedie.
\newblock An index which distinguishes among some important exponential families.
\newblock {\em Statistics: Applications and New Directions}, 579--604, 1984.

\bibitem{schepsmeier2016}
U.~Schepsmeier and C.~Czado.
\newblock Vine copula-based bias correction of climate model outputs.
\newblock {\em Environmental and Ecological Statistics}, 23(4):573--597, 2016.

\bibitem{sklar1959}
M.~Sklar.
\newblock Fonctions de r{\'e}partition {\`a} n dimensions et leurs marges.
\newblock {\em Publications de l'Institut de Statistique de l'Universit{\'e} de Paris}, 8:229--231, 1959.

\bibitem{joe2014}
H.~Joe.
\newblock {\em Dependence Modeling with Copulas}.
\newblock CRC Press, 2014.

\bibitem{genest2007}
C.~Genest and J.~Ne{\v{s}}lehov{\'a}.
\newblock On copulas for multivariate discrete data.
\newblock {\em Astin Bulletin}, 37(2):475--515, 2007.

\bibitem{barrientos2012}
A.~F. Barrientos, A.~Jara, and F.~A. Quintana.
\newblock On the support of MacEachern's dependent Dirichlet processes.
\newblock {\em Bayesian Analysis}, 7(2):277--310, 2012.

\bibitem{lavine1992}
M.~Lavine.
\newblock Some aspects of Polya tree distributions for statistical modelling.
\newblock {\em The Annals of Statistics}, 20(3):1222--1235, 1992.

\bibitem{hanson2002}
T.~E. Hanson and W.~O. Johnson.
\newblock Modeling nominal data using Polya trees.
\newblock {\em Journal of Nonparametric Statistics}, 14(3):285--299, 2002.

\bibitem{hanson2006}
T.~E. Hanson.
\newblock Inference for mixtures of Polya trees.
\newblock {\em The Annals of Statistics}, 34(3):1248--1286, 2006.

\bibitem{ishwaran2001}
H.~Ishwaran and L.~F. James.
\newblock Gibbs sampling for Dirichlet process mixture models with a search for the number of clusters.
\newblock {\em Journal of the American Statistical Association}, 96(455):1061--1073, 2001.

\bibitem{chipman2010}
H.~A. Chipman, E.~I. George, and R.~E. McCulloch.
\newblock BART: Bayesian additive regression trees.
\newblock {\em The Annals of Applied Statistics}, 4(1):266--298, 2010.

\bibitem{linero2018}
A.~R. Linero.
\newblock Bayesian regression trees for high-dimensional prediction and variable selection.
\newblock {\em Journal of the American Statistical Association}, 113(522):626--636, 2018.

\bibitem{bleich2014}
J.~Bleich, A.~Kapelner, E.~I. George, and S.~T. Jensen.
\newblock Variable selection workflows for Bayesian additive regression trees.
\newblock {\em The Annals of Applied Statistics}, 8(3):1750--1781, 2014.

\bibitem{ishwaranrao2005}
H.~Ishwaran and J.~S. Rao.
\newblock Spike and slab variable selection: frequentist and Bayesian strategies.
\newblock {\em The Annals of Statistics}, 33(2):730--773, 2005.

\end{thebibliography}

\end{document}